\textit{
Bank Indonesia's public Project Garuda experiments focus on wholesale central
bank digital currency. This study asks how a testable retail Digital Rupiah
prototype can represent two-tier distribution, wallet policies, and
institutional oversight on a permissioned ledger. It uses Hyperledger Fabric
2.5, Go chaincode and backend services, CouchDB, and PostgreSQL.
}

\textit{
Raw identity data remains in PostgreSQL; the ledger stores KYC status, risk,
expiry, and document hashes. This separation does not provide transaction
privacy on the shared channel. Banks or payment service providers supply KYC
decisions. The chaincode derives BASIC, STANDARD, and MERCHANT wallet tiers and
validates transaction, daily, monthly, and balance limits. Backend ownership
and chaincode MSP checks enforce user and organization boundaries. These
controls do not demonstrate complete AML/CFT compliance or four-eyes approval.
}

\textit{
Automated tests and exact-diagnostic negative paths evaluate authorization and
policy enforcement. Hyperledger Caliper evaluates a fixed retail population
under 10, 30, and 60 TPS offered loads, with a clean ledger for each profile,
deterministic seeding, repeated peak rounds, and an aggregate-overspend
scenario. At a 30 TPS offered load, successful throughput reaches 29.5 TPS
without failures; at 60 TPS, it ranges from 41.2 to 47.0 TPS. All seven negative
cases are rejected for the expected reason. In the overspend scenario, 4 of 200
transfers commit and the final balance remains nonnegative. These results form a
repeatable local baseline, not evidence of cross-machine reproducibility or
production capacity. The prototype is an auditable artifact for retail Digital
Rupiah policy experiments within the stated limitations.
}

\vspace{0.5cm}

\noindent\textbf{Keywords}: Digital Rupiah, retail CBDC, Hyperledger Fabric,
risk-based KYC, wallet limits
